\subsubsection{Prim (priority queue)}

\paragraph{Idea intuitiva.}
Otra forma de calcular el MST: empezar desde un nodo y, en cada paso, aniadir la arista de menor peso que conecta el conjunto construido con el resto. Usando priority queue, es $O((V + E) \log V)$. La demostracion: cada vez que se aniade una arista, es la minima cruzando el ``corte'' entre el conjunto actual y el resto (Cut Property); eso preserva optimalidad.

\paragraph{Propiedades clave (invariantes).}
\begin{itemize}\setlength\itemsep{0pt}
	\item Similar a Dijkstra pero sin acumular distancias, solo eligiendo minimo por arista.
	\item Funciona mejor que Kruskal en \textbf{grafos densos} ($E \approx V^2$).
	\item Si el grafo es disconexo, genera el bosque (solo llega a los alcanzables).
	\item Cada nodo se inserta al heap una vez; cada arista se procesa una vez.
\end{itemize}

\paragraph{Codigo clave.}
Prim clasico con heap:
\begin{verbatim}
priority_queue<pair<int, int>, vector<...>, greater<...>> pq;
vector<bool> used(n, false);
vector<int> minEdge(n, INF);
minEdge[0] = 0;
pq.push({0, 0});
while (!pq.empty()) {
    auto [w, v] = pq.top(); pq.pop();
    if (used[v]) continue;
    used[v] = true;
    mst += w;
    for (auto [peso, u] : adj[v])
        if (!used[u] && peso < minEdge[u]) {
            minEdge[u] = peso;
            pq.push({peso, u});
        }
}
\end{verbatim}

\paragraph{Complejidad.}
\begin{itemize}\setlength\itemsep{0pt}
	\item $O((V + E) \log V)$.
	\item Con matriz de adyacencia: $O(V^2)$ sin heap.
\end{itemize}

\paragraph{Donde tocar para modificar.}
\begin{itemize}\setlength\itemsep{0pt}
	\item \textbf{Prim con matriz de adyacencia}: $O(V^2)$ sin heap, util en grafos densos pequenos.
	\item \textbf{Reconstruir el MST}: aniadir \texttt{parent[]}.
	\item \textbf{Prim en streaming}: variante para grafos que llegan incrementalmente.
\end{itemize}

\paragraph{Trampas y errores frecuentes.}
\begin{itemize}\setlength\itemsep{0pt}
	\item El \texttt{if (used[v]) continue;} es crucial: si el nodo ya fue procesado, saltar.
	\item En la version matriz, elije el minimo por linea (no heap).
	\item Si el grafo es disconexo, aniadir un loop externo sobre todos los nodos no usados.
\end{itemize}

\paragraph{Codigo.}
Ver \texttt{cpp.tex}, seccion \textit{Prim (priority queue)}.
